\documentclass[UTF8,a4paper,11pt]{ctexart}

\usepackage[a4paper,left=2.2cm,right=2.2cm,top=2.1cm,bottom=2.1cm]{geometry}
\usepackage{amsmath,amssymb,bm}
\usepackage{booktabs,tabularx,array,longtable,multirow}
\usepackage{enumitem}
\usepackage{xcolor}
\usepackage{hyperref}
\usepackage{xurl}
\usepackage{fancyhdr}
\usepackage{titlesec}
\usepackage{tcolorbox}
\usepackage{parskip}

\definecolor{navy}{RGB}{29,59,92}
\definecolor{lightblue}{RGB}{235,244,252}
\definecolor{lightgreen}{RGB}{237,248,241}
\definecolor{lightorange}{RGB}{255,246,232}
\definecolor{graytext}{RGB}{75,75,75}
\hypersetup{colorlinks=true,linkcolor=navy,urlcolor=blue,citecolor=navy}
\setlist{nosep,leftmargin=2em}
\setlist[itemize]{label=\textbullet}
\setlength{\parindent}{2em}
\setlength{\emergencystretch}{5em}
\setlength{\tabcolsep}{5pt}
\renewcommand{\arraystretch}{1.25}
\pagestyle{fancy}
\fancyhf{}
\fancyhead[L]{2.5D 芯粒布局：状态化宏动作与算法组合}
\fancyhead[R]{实验方案底稿}
\fancyfoot[C]{\thepage}
\setlength{\headheight}{14pt}
\titleformat{\section}{\Large\bfseries\color{navy}}{\thesection}{0.6em}{}
\titleformat{\subsection}{\large\bfseries\color{navy}}{\thesubsection}{0.6em}{}

\newcommand{\code}[1]{\texttt{#1}}
\newcommand{\important}[1]{\textbf{#1}}
\newcommand{\urlbreak}[1]{\url{#1}}

\begin{document}

\begin{titlepage}
\centering
\vspace*{1.2cm}
{\LARGE\bfseries 面向2.5D系统集成的芯粒布局方法\par}
\vspace{0.35cm}
{\Huge\bfseries 状态化宏动作与 CNN--Contextual Bandit 实验设计方案\par}
\vspace{0.8cm}
{\Large ——基于状态表征、宏动作选择与统一物理反馈的多算法布局优化\par}
\vspace{1.2cm}
\begin{tcolorbox}[colback=lightblue,colframe=navy,width=0.92\textwidth,boxrule=0.7pt]
\centering
\large
本方案的核心不是把 SA、强化学习和解析法串成流水线，而是学习一个\important{状态条件化的算法组合策略}：\par
\vspace{0.2cm}
\(\text{布局状态}\;\longrightarrow\;\text{CNN 上下文表征}\;\longrightarrow\;\text{选择算法/配置/预算}\;\longrightarrow\;\text{物理评估反馈}\)。
\end{tcolorbox}
\vfill
\begin{flushleft}
\textbf{文档定位：}正式实验前的实验设计、接口定义、消融方案和开源实现清单。\\
\textbf{当前方法版本：}当前已实现 fourmap v2 状态投影；拟采用四通道 2D CNN + 有限配置库 + Contextual Bandit（CNN 尚未训练）；后续按需扩展到 CNN+GRU、ConvLSTM 和低分辨率局部注意力。\\
\textbf{重要边界：}本文档中的性能结论均为待验证的实验假设，不将已有项目的报告数值当作本项目结果。
\end{flushleft}
\vspace{0.8cm}
\today
\end{titlepage}

\tableofcontents
\newpage

\section{研究问题与总体判断}

\subsection{研究张力}

2.5D 芯粒布局同时受到互连线长、峰值温度、面积利用率、芯粒不重叠和边界合法性等约束。不同算法的优势并不对所有状态保持不变：模拟退火擅长在局部邻域中持续探索，解析法通常具有较强的目标函数驱动和较快的收敛，强化学习能够利用反复交互形成策略，但它们对初始布局、热点分布、连接密度和可行域形状的敏感性不同。

因此，真正需要研究的问题不是“如何让三个算法依次工作”，而是：

\begin{quote}
在给定布局状态和统一物理评估预算下，能否从状态的空间分布与历史变化中预测各个算法配置的短期收益，并在线选择最合适的宏动作？
\end{quote}

这里的三类 agent 应被定义为\important{三种算法 worker}：SA-Agent、Analytical-Agent 和 RL-Agent。它们分别产生候选布局；上层 Contextual Bandit/Coordinator 学习在什么状态下调用哪一种 worker、使用什么配置、wiremask 和宏动作预算。它们不直接替代物理规律，也不直接生成未经评估的可信布局坐标；物理评估器仍然是最终裁决者。

\subsection{贡献增量的准确表述}

本项目相对于“固定使用一个算法”或“把多个算法首尾连接”的具体增量应写成以下三点。

\begin{enumerate}
    \item \textbf{状态条件化的算法选择：}将当前布局编码为包含连接优先级、热代理、HPWL 和合法性诊断信息的空间状态张量，Contextual Bandit 根据该状态选择 \((\text{algorithm},\text{configuration},\text{budget})\) 宏动作。
    \item \textbf{统一接口下的可替换布局算子：}SA、解析法和 RL worker 均接收同一种布局状态、输出同一种候选布局，并由同一个 evaluator 返回可比指标；它们是可竞争的动作，不是预先规定的阶段链路。
    \item \textbf{把“算法偏好”变成可测量的状态转移：}除了记录最终结果，还记录算法作用前后的状态差分、空间影响图和资源成本，从而回答“某算法在何种布局结构上有效、改善了哪个指标、代价是什么”。
\end{enumerate}

\begin{tcolorbox}[colback=lightgreen,colframe=green!45!black,title=论文主线]
\textbf{空间状态表征}定义了决策上下文；\textbf{Contextual Bandit}定义了算法选择；\textbf{布局 worker}提供候选动作；\textbf{统一物理评估器}提供奖励与约束；\textbf{状态流数据}使决策器学习算法的适用条件。整个系统的创新点位于“状态--算法--结果”之间的可学习映射，而不在于简单增加 agent 数量。
\end{tcolorbox}

\section{总体架构：闭环决策而非流水线拼接}

\subsection{系统流程}

每个决策时刻 \(t\) 的流程定义为：

\begin{equation}
\begin{aligned}
X_t &\xrightarrow{\text{CNN encoder}} z_t
\xrightarrow{\text{Contextual Bandit}} a_t=(k_t,c_t,H_t,M_t)
\\[-0.2em]
&\xrightarrow{\text{worker }k_t}
Y_{t+1}
\xrightarrow{\text{Physics Evaluator}}
(X_{t+1},m_{t+1},r_t).
\end{aligned}
\end{equation}

其中 \(X_t\) 是当前布局状态，\(k_t\) 是算法类别，\(c_t\) 是算法配置，\(H_t\) 是本次宏动作的预算，\(Y_{t+1}\) 是候选布局，\(m_{t+1}\) 是合法性和物理指标，\(r_t\) 是控制器根据目标函数计算的奖励。只有当候选布局满足硬合法性约束或经过确定性修复后，才允许进入下一次决策。

\subsection{决策单位：微步、宏步与完整优化过程}

本方案将 Contextual Bandit 的“一步”定义为一次\important{宏动作（macro-action）}，而不是 SA 的一次随机移动、RL 的一次环境动作，也不是让某个算法一直运行到收敛。三个层级必须严格区分：

\begin{table}[htbp]
\centering
\small
\begin{tabularx}{\textwidth}{p{0.18\textwidth}p{0.34\textwidth}X}
\toprule
层级 & 定义 & 记录与控制方式 \\
\midrule
微步（micro-step） & SA 的一次邻域移动、解析法的一次更新或 RL 的一个环境动作 & 由 worker 内部执行，可用于调试和重放，不由上层 Bandit 逐次选择 \\
宏步（macro-step） & 一个 worker 在固定 evaluator/时间预算下连续执行若干微步后，返回一个候选布局 & Bandit 选择一次，记录起点状态、动作、终点状态、指标差分和资源成本 \\
episode & 从初始布局到停止条件的完整优化过程 & 记录宏动作序列、累计收益、总预算和最终布局 \\
\bottomrule
\end{tabularx}
\end{table}

因此动作定义为
\begin{equation}
a_t=(k_t,c_t,H_t,M_t),
\end{equation}
其中 \(k_t\) 是算法类别，\(c_t\) 是离散配置，\(H_t\) 是 evaluator 或迭代预算，\(M_t\) 是本次动作使用的 wiremask 及其约束。宏步在达到预算、worker 主动停止、超时或产生不可修复的非法布局时结束。该定义既允许 agent 自主选择算法和参数，又避免让控制器承担三个算法内部循环的全部细节。

\begin{tcolorbox}[colback=lightorange,colframe=orange!70!black,title=实验边界]
第一版只学习有限的 \((\text{算法},\text{配置},\text{预算})\) arm；连续超参数优化不与算法选择同时作为首个研究问题。只有当有限 arm 版本证明状态能够区分算法适用条件后，才引入分层或连续参数 Bandit，并将参数搜索成本计入 reward。
\end{tcolorbox}

\begin{tcolorbox}[colback=lightorange,colframe=orange!70!black,title=与原多阶段 Agent 方案的关系]
原方案中的 order、thermal、geometry 不再作为三个固定顺序的 agent。它们被改写为 \textbf{可学习的诊断通道/动作约束}：连接优先级、热代理热点和几何冲突分别进入状态张量与 wiremask；最终实际执行的是一个算法 worker。若实验发现某类状态稳定需要先做热修复，再可把“热修复宏动作”加入动作库，但这个顺序应由对照实验支持，而不是在系统结构中预先假定。
\end{tcolorbox}

\subsection{统一状态表示}

对于给定的 chiplet 集合和板边界，将布局栅格化为 $H\times W$，默认 $256\times256$。当前 fourmap v2 的唯一基础输入为
\begin{equation}
X_t=[M_{\rm priority},M_{\rm thermal\_proxy},M_{\rm hpwl},M_{\rm legality}]
\in\mathbb{R}^{4\times H\times W}.
\end{equation}
全局标量 $v_t$（精确合法性、重叠、越界、HPWL、热代理、预算和最近动作）单独融合，不计为空间通道。占用、功耗、连接图和 wiremask 可作为诊断或后续扩展，均不隐式加入当前四通道。历史三联图保留原版本；新投影不覆盖旧图、评估器或已记录奖励。具体定义见后文。

\paragraph{wiremask 的位置。} 本设计中 wiremask 不是额外的“第四个优化阶段”，也不是只在最后筛选区域。它有两个作用：

\begin{enumerate}
    \item 作为由四通道导出的可选门控，使 Bandit 看到“当前最需要被改变的空间区域”；
    \item 作为 worker 的动作约束，将候选移动、局部搜索或 RL 动作限制在允许区域/允许方向内，减少无效探索。
\end{enumerate}

第一版不要求 wiremask 具有端到端可学习的生成器。建议先用确定性规则生成 baseline mask，再把 mask 的阈值、膨胀半径和指标权重加入配置库，作为可控消融因素。

建议同时记录 hard mask、soft mask 和 mask version。hard mask 用于限制候选移动区域、方向和移动半径；soft mask 作为可选 CNN/注意力门控的连续权重；mask version 用于保证不同实验不会隐式使用不同的候选区域。只有在确定性 mask 的离线结果证明有收益后，才让诊断头预测软 mask。

\section{分阶段实验设计}

\subsection{阶段 0：基准数据、统一格式与物理评估器}

\paragraph{目标。}在任何 Bandit 训练前，先建立唯一的 canonical layout representation，避免三个算法各自使用不同的坐标、线长定义或热评价器，导致结果不可比。

建议的数据结构如下：

\begin{table}[htbp]
\centering
\small
\begin{tabularx}{\textwidth}{>{\raggedright\arraybackslash}p{0.22\textwidth}X}
\toprule
字段 & 内容 \\
\midrule
几何 & chiplet ID、宽高、方向、坐标、die 边界、禁止区域、间距约束 \\
连接 & net 列表、pin offset、net weight、芯粒连接矩阵 \\
热 & 功耗、材料/层配置、散热边界条件、热模型版本 \\
状态 & 栅格化通道、全局指标、当前预算、历史动作和随机种子 \\
输出 & 候选布局、合法性、HPWL、路由线长、峰值/平均温度、运行时间、调用次数 \\
\bottomrule
\end{tabularx}
\end{table}

\paragraph{阶段 0 的验收。} 对每个基准 case，给定同一个布局文件，三种 worker 的 evaluator 输出应在容差内一致；至少测试边界、重叠、旋转、pin offset、HPWL 和功耗读取。若同一布局在三个 worker 中得到不同 HPWL，不进入 Bandit 训练，先修正 adapter。

\subsection{阶段 1：三类算法 Worker 与有限配置库}

第一版不要同时进行连续超参数优化。将算法和关键参数离散成有限配置库，使 Bandit 先解决一个定义清晰的有限动作选择问题：

\begin{equation}
\mathcal{A}=\{(k,c): k\in\{\mathrm{SA},\mathrm{AN},\mathrm{RL}\},c\in\mathcal{C}_k\}.
\end{equation}

其中 AN 表示 analytical placement。每个 arm 在固定的 evaluator 预算 \(H\) 下运行一次，返回宏观结果，不把内部每一个随机移动都交给 Bandit。

\begin{table}[htbp]
\centering
\small
\begin{tabularx}{\textwidth}{p{0.15\textwidth}p{0.24\textwidth}X}
\toprule
Worker & 第一版配置字段 & 推荐的第一批离散配置 \\
\midrule
SA & 初温、降温率、每温度迭代数、移动半径、wire/thermal 权重、局部/全局搜索比例 & 低/中/高三档温度；快/慢两档降温；小/大移动半径；线长主导/热主导/均衡三种权重 \\
AN & wire 权重、thermal 权重、密度权重、学习率、温度阈值、热感知开关 & WL-driven、Thermal-aware、两种热权重、两种密度权重，组成不超过 8--12 个配置 \\
RL & checkpoint、动作步长、熵系数、折扣因子、候选动作半径、热模型版本 & 预训练 policy 的小/中/大动作半径；低/中探索；固定 policy 与短 horizon fine-tune 两档 \\
\bottomrule
\end{tabularx}
\end{table}

\paragraph{为什么先用有限配置库。} 当前研究问题是“状态是否能指导算法选择”，不是“连续超参数优化是否有效”。有限 arm 可以计算选择 regret、top-$k$ 命中率和每个状态簇的算法偏好。如果一开始同时做算法选择和连续参数优化，性能提升无法归因，且会把 Bandit、Bayesian optimization 和 worker 内部训练混在一起。

\paragraph{连续参数的后续版本。} 第二阶段可采用分层动作：第一层选择算法 \(k\)，第二层在该算法的可信参数范围内用离散化、Gaussian Process Bandit 或 Thompson Sampling 优化连续参数 \(\theta\)。此时必须将参数搜索成本纳入 reward，且保留有限库版本作为基线。

\subsection{阶段 2：算法轨迹、状态转移与训练样本}

\paragraph{轨迹与状态流的区别。} 三种数据对象必须分开命名：

\begin{table}[htbp]
\centering
\small
\begin{tabularx}{\textwidth}{p{0.2\textwidth}X}
\toprule
对象 & 定义与用途 \\
\midrule
算法轨迹 & worker 内部的完整事件序列，如 SA 的每次移动/接受概率，RL 的每步 observation/action，解析法的每次迭代。它保留细节，适合调试和重放。 \\
宏观状态转移样本 & 一次宏动作前后记录 \((X_t,a_t,X_{t+1},r_t,c_t)\)，其中 \(a_t\) 是算法、配置和预算，\(c_t\) 是时间/evaluator 调用等成本。它是 Bandit 的基本训练单元。 \\
状态流 & 将连续快照抽象为序列 \(X_t\rightarrow X_{t+1}\rightarrow\cdots\)，重点学习某种算法改变空间分布和指标的规律，而不关心每一次内部事件。它是 3D CNN/ConvLSTM 的输入。 \\
\bottomrule
\end{tabularx}
\end{table}

也就是说，轨迹是“发生过哪些内部事件”，状态流是“状态如何变化”。第一版只需保存宏观状态转移样本；第二版再使用固定时间间隔的连续状态快照构造状态流。

\paragraph{数据采集协议。}

\begin{enumerate}
    \item 按 benchmark instance 划分训练、验证和测试，不能将同一个 case 的不同中间状态随机分到不同集合，以免泄露布局结构。
    \item 对每个初始布局，让所有 arm 在同一 evaluator 预算、相同随机种子集合和相同输入状态下独立运行。
    \item 每个宏动作保存起点快照、终点快照、增量 mask、指标变化、运行时间、evaluator 调用次数和 worker 配置。
    \item 早期必须执行一轮全臂采样，避免 Bandit 只看到先验偏好的算法；后续再使用 epsilon-greedy 或 Thompson Sampling 探索。
\end{enumerate}

\subsection{阶段 3：2D CNN + Contextual Bandit}

\paragraph{模型。} 第一版使用共享 CNN encoder：

\begin{equation}
z_t=f_\phi(X_t),\qquad
\hat{q}_t(a)=g_\psi([z_t,e_a]),
\end{equation}

其中 \(e_a\) 是算法/配置的 embedding，\(\hat{q}_t(a)\) 预测执行 arm \(a\) 后的短期收益。共享 encoder 的意义是让不同算法基于同一个物理状态空间竞争，避免训练三个互不对齐的 CNN。

推荐顺序：

\begin{enumerate}
    \item 先离线训练 CNN 回归器，预测固定预算 \(H\) 下的 \(\Delta\)HPWL、热代理变化、精确合法性和运行成本；有独立高保真标签时再预测 \(\Delta\)Tpeak；
    \item 再将 CNN 表征接入 LinUCB 或 linear Thompson Sampling；
    \item 最后做 online update，使用真实 evaluator 返回的 reward 修正预测。
\end{enumerate}

\paragraph{奖励。} 不直接把单一的“总分”当作唯一事实，而是记录原始指标并由实验配置产生目标化奖励：

\begin{equation}
r_t = -\lambda_w\Delta \widetilde{WL}_t
       -\lambda_T\Delta \widetilde{T}_t
       -\lambda_A\Delta \widetilde{A}_t
       -\lambda_L\,\mathbb{1}[\text{illegal}]
       -\lambda_C\widetilde{C}_t.
\end{equation}

其中当前热项为 thermal proxy，有独立温度标签时才使用温度项，二者不混用。\(\Delta\) 表示相对当前状态的变化，\(\widetilde{C}_t\) 是归一化的运行时间或 evaluator 调用成本。所有 arm 使用相同的成本口径；否则 Bandit 可能只学会选择快而不一定有效的 worker。

\paragraph{Bandit 的评价。} 第一版不只报告最终布局分数，还报告：

\begin{itemize}
    \item arm 选择 regret：与事后知道最优 arm 的 oracle 比较；
    \item top-1/top-3 选择准确率：预测收益最高的 arm 是否真正有效；
    \item reward 预测误差和校准误差；
    \item 在相同 evaluator budget 下的最终 HPWL、峰值温度、合法率和 wall-clock time。
\end{itemize}

\subsection{阶段 4：3D CNN/ConvLSTM 学习状态流}

第二版不是“把 2D CNN 换成更大的网络”，而是改变学习目标：预测某个算法宏动作造成的状态变化。推荐先用 CNN/ECA + GRU 处理压缩后的历史特征；只有当历史本身有增益且空间位置仍然重要时，再升级为小型 ConvLSTM。全局 Transformer 不作为默认第二版。

给定长度为 \(L\) 的快照序列：

\begin{equation}
\mathcal{S}_{t-L+1:t}=[X_{t-L+1},\ldots,X_t],
\end{equation}

3D CNN 或 ConvLSTM 估计：

\begin{equation}
\widehat{F}_{k,c}(\mathcal{S}_{t-L+1:t})
\rightarrow
(\widehat{X}_{t+1},\widehat{\Delta m}_{t+1},\widehat{\Delta C}_{t+1}).
\end{equation}

输入的时间维度刻画“最近几步变化方向”，空间维度刻画热点/拥挤区域。Bandit 的上下文变为 \(\widehat{F}_{k,c}\) 对不同候选动作的预测，而不再只依赖当前静态状态。

\paragraph{第二版的必要约束。} 只有当第一版证明静态状态不能区分不同算法的适用条件时，才进入 3D/ConvLSTM。否则 3D 模型增加数据量和训练复杂度，却无法证明状态流带来了决策收益。建议使用以下消融：

\begin{enumerate}
    \item 2D CNN 静态上下文；
    \item 2D CNN + 全局历史标量；
    \item 3D CNN/ConvLSTM + 状态差分通道；
    \item 3D 模型去掉 wiremask 或去掉 thermal 通道。
\end{enumerate}

\section{状态化扩展：诊断、记忆与流程级收益}

\subsection{从静态上下文到三个可验证的预测任务}

为了避免把 CNN、注意力和 Bandit 拼成并列模块，建议把控制器拆成三个具有不同实验含义的输出头：

\begin{enumerate}
    \item \textbf{诊断头 \(D_t\)：}预测当前最值得干预的区域、指标和异常类型，例如线长拥挤、热热点或合法性风险；
    \item \textbf{转移头 \(F_t(a)\)：}给定候选宏动作 \(a\)，预测它对空间 map 和全局指标的变化；
    \item \textbf{选择头 \(Q_t(a)\)：}综合预测收益、执行成本、不确定性和后续价值，决定是否选择该 arm。
\end{enumerate}

该划分使“哪里有问题”“执行后会怎样”和“现在是否值得执行”成为三个可分别验证的假设。注意力权重只能作为诊断线索，不能直接等同于物理因果解释；必须用 evaluator 的实际增量 map 和指标差分进行校验。

\subsection{推荐的轻量网络路线}

不建议在 \(256\times256\) 的像素级 map 上直接使用全局 Transformer。推荐按成本由低到高推进：

\begin{enumerate}
    \item \textbf{第一版：CNN/ECA + Contextual Bandit。}共享 CNN 提取局部结构，ECA 或 CBAM 风格的轻量通道/空间门控突出重要 map 通道，LinUCB 或 linear Thompson Sampling 选择有限 arm；
    \item \textbf{第二版：CNN/ECA + GRU。}将最近 \(L\) 个宏步压缩为全局特征和 memory packet，先检验“历史变化方向”是否有增益；
    \item \textbf{第三版：小型 ConvLSTM。}只有当 GRU 证明空间位置对状态转移有额外贡献时，才保留空间维度建模；
    \item \textbf{第四版：低分辨率局部注意力。}将 map 下采样为 \(16\times16\) 或 \(32\times32\) patch token，使用局部窗口、shifted-window 或 chiplet/热点 token 进行跨图和跨时间交互。
\end{enumerate}

这个顺序的理由不是网络越大越好，而是每一步都对应一个明确的可证伪问题：空间信息是否有用、历史是否有用、空间化历史是否有用、局部注意力是否进一步改善决策。

\subsection{状态输入与结构化 memory packet}

当前状态仍为 $X_t\in\mathbb{R}^{4\times H\times W}$，$v_t$ 和阶段信息保存在结构化旁路。后续可将最近 $L$ 步的四通道状态与动作组成历史窗口，比较 2D CNN + history、CNN+GRU 和 ConvLSTM；这些均为未来实验，不改变当前基础通道定义。

内循环不使用自然语言 prompt 传递状态，而使用可序列化的结构化记忆（以下为接口示例，非实验结果）：

\begin{verbatim}
{
  "step": 7,
  "phase": "repair_or_refine",
  "focus_region": [[x0, y0, x1, y1]],
  "focus_metric": "hpwl",
  "mask_version": "wire-v2",
  "last_action": {"worker": "sa", "config": "sa-wire", "budget": 4},
  "predicted_delta": {"hpwl": null, "thermal_proxy": null},
  "observed_delta": {"hpwl": null, "thermal_proxy": null},
  "uncertainty": null,
  "remaining_budget": 96
}
\end{verbatim}

GRU/ConvLSTM 读取该信息的张量化版本。若后续使用 LLM，它只承担外层实验配置、异常审计或可读解释，不参与每个 inner-loop 的算法选择，从而保持控制协议可复现。

\subsection{wiremask 的定义与使用}

wiremask 由基础四通道导出，作为可选门控和动作约束：

\begin{enumerate}
    \item \textbf{可选门控：}由 HPWL 贡献、net bounding box、连接拥挤度、热图和合法性膨胀规则生成，告诉编码器哪些区域最值得关注；
    \item \textbf{动作约束：}限制 SA 邻域、解析更新或 RL 动作的候选区域、方向、移动半径和可行域，减少无效 evaluator 调用；
    \item \textbf{状态差分：}记录动作前后的 HPWL 图、导出 mask 及其变化，用于训练转移头和评估 mask 是否真正改变了布局；
    \item \textbf{可控消融：}第一版使用确定性 mask，阈值、膨胀半径和通道权重进入有限配置库；只有在离线结果证明有收益后，才让诊断头预测软 mask。
\end{enumerate}

建议同时记录 hard mask、soft mask 和 mask version，避免不同实验隐式使用不同的候选区域。

\subsection{从单步 reward 到流程级价值}

纯 Contextual Bandit 只估计当前宏步的短期收益。为体现“这一步对整个流程的作用”，按以下三层推进：

\begin{equation}
r_t^{\rm local}=J(X_t)-J(X_{t+1})-\lambda C_t-\rho I_t,
\end{equation}
其中 \(J\) 是归一化目标，\(C_t\) 是时间或 evaluator 调用成本，\(I_t\) 是非法或修复惩罚。

当短期基线稳定后，使用 \(K\)-step return：
\begin{equation}
R_t^{(K)}=\sum_{i=0}^{K-1}\gamma^i r_{t+i}+\gamma^K V(X_{t+K}).
\end{equation}

进一步学习动作条件的 flow critic：
\begin{equation}
F_\theta(X_t,a_t)\rightarrow
(\widehat{\Delta M},\widehat{\Delta v},\widehat{cost},\widehat{uncertainty}).
\end{equation}

其中 \(\widehat{\Delta M}\) 描述 wire/thermal/legal map 的变化，\(\widehat{\Delta v}\) 描述全局指标变化。这样控制器比较的是一个动作在后续流程中的预期作用，而不只是下一步分数。


\subsection{轨迹、宏观转移与状态流}

三类数据对象必须分开保存：

\begin{table}[htbp]
\centering
\small
\begin{tabularx}{\textwidth}{p{0.2\textwidth}X}
\toprule
对象 & 含义与用途 \\
\midrule
算法轨迹 & worker 内部的完整事件序列，如 SA 的每次移动/接受概率、RL 的每步 observation/action 和解析法的每次更新；用于调试和重放。 \\
宏观状态转移 & \((X_t,a_t,X_{t+1},\Delta X,\Delta m,c_t)\)；是第一版 CNN + Bandit 的基本训练单元。 \\
状态流 & 固定宏步时间尺度上的 \(X_t\rightarrow X_{t+1}\rightarrow\cdots\)，强调状态如何变化；是 GRU/ConvLSTM/时序注意力的输入。 \\
\bottomrule
\end{tabularx}
\end{table}

第一版只需完整保存宏观状态转移；第二版再按固定窗口构造状态流。数据集必须按 benchmark instance 划分，不能将同一个 case 的不同中间状态随机拆到训练集和测试集。

\section{物理评估器设计}

\subsection{评估器不是黑盒奖励函数}

Physics Evaluator 应返回原始物理量、合法性、计算成本和模型版本；奖励由 controller 计算。推荐接口：

\begin{verbatim}
evaluate(layout, fidelity, case_config) -> {
    legal, overlap, boundary_violation,
    hpwl, routed_wirelength,
    peak_temperature, mean_temperature,
    area_utilization, runtime, evaluator_calls,
    fidelity, evaluator_version
}
\end{verbatim}

不能只返回一个 reward，否则无法判断是线长改善、温度恶化还是成本增加导致的结果。

\subsection{三个物理子评估器}

\begin{enumerate}
    \item \textbf{GeometryEvaluator：}以精确几何检查边界、重叠、旋转、间距和禁止区域；栅格仅供诊断。它是合法性的唯一裁判，必须确定性执行，并在每次宏动作后调用。
    \item \textbf{WirelengthEvaluator：}内环使用带 pin offset 和 net weight 的 HPWL 或快速中心距离近似；高保真阶段可调用 TAP-2.5D 的路由/MILP 接口。报告 HPWL 与 routed wirelength 两个不同字段。
    \item \textbf{ThermalEvaluator：}内环使用紧凑解析热模型或快速热模型；定期抽样和最终验证使用 HotSpot wrapper。快速模型只能提供低成本反馈，不能被表述为最终热真实性本身。
\end{enumerate}

\subsection{多保真度与公平性}

建议定义三种 fidelity：

\begin{table}[htbp]
\centering
\small
\begin{tabularx}{\textwidth}{p{0.16\textwidth}p{0.3\textwidth}X}
\toprule
等级 & 组成 & 使用阶段 \\
\midrule
Fast & 几何合法性 + 增量 HPWL + 紧凑热模型 & Bandit 训练内环、worker 的大多数中间步 \\
Medium & 几何合法性 + HPWL + 更精细热模型 & 数据采样和候选动作比较 \\
Exact & 几何合法性 + 路由线长 + HotSpot 热评价 & 验证集、最终布局、跨模型校准 \\
\bottomrule
\end{tabularx}
\end{table}

公平性要求：

\begin{itemize}
    \item 同一次比较中，所有算法使用相同 fidelity、最大 wall-clock 或 evaluator call budget；
    \item 训练 RL policy 的成本和部署时 inference 成本分开记账；
    \item 快速热模型和 HotSpot 的偏差用独立校准集报告，不在测试集上反复调参；
    \item 所有结果记录 evaluator 版本、输入单位、随机种子和源码 commit hash。
\end{itemize}

\section{对照实验与消融矩阵}

\subsection{主实验}

\begin{table}[htbp]
\centering
\small
\begin{tabularx}{\textwidth}{p{0.24\textwidth}X}
\toprule
实验组 & 研究问题 \\
\midrule
固定 SA & 单一经典搜索器的性能上限与成本 \\
固定 analytical & 目标函数驱动方法的性能与收敛特征 \\
固定 RL & 学习型策略单独使用的效果与泛化 \\
随机选择 arm & 多算法组合但无状态决策是否有效 \\
全局指标 Bandit & 只有标量上下文时，空间状态是否必要 \\
2D CNN + Bandit & 主方法：静态空间状态能否改善算法选择 \\
2D CNN + wiremask & 派生 mask 作为门控和动作约束是否进一步改善选择与局部修复 \\
CNN + GRU & 最近宏步的历史是否提升收益预测与在线选择 \\
3D/ConvLSTM + Bandit & 状态流是否带来超过静态状态的增益 \\
ConvLSTM/局部注意力 + flow critic & 动作条件的状态转移预测是否改善流程级质量--成本折中 \\
\bottomrule
\end{tabularx}
\end{table}

\subsection{必须控制的变量}

比较时固定：benchmark split、初始布局、芯粒尺寸和功耗、目标权重、evaluator fidelity、总预算、随机种子数和最终验证流程。若某算法需要预训练，应分别报告“包含训练成本的端到端结果”和“已训练 worker 的部署结果”，避免将训练资源隐藏在基线之外。

\subsection{三对四通道与优先级消融协议}

在同一 canonical 状态集合上重新投影，比较 v2 的前三通道与完整四通道，唯一输入差异为 legality；历史三联图仅作归档对照，避免旧 priority 覆盖规则混淆第四通道效应。优先级消融固定其余三通道，比较完整 $q_i$、删除 priority、仅保留占用（芯粒 0、背景 $-1$）以及状态内打乱芯粒 $q_i$ 的版本；打乱保留值的多重集与并列，使用固定种子。

各组固定相同状态 ID、case holdout、全臂标签、奖励、evaluator fidelity、候选配置、随机种子集合和评价预算；同一 case 的全部中间状态及同源变体不得跨训练/测试集。归一化和调参只用训练/验证 case，匹配训练与调参预算，单列参数量与投影成本。当前 129 个状态不能当作 129 个独立 case。报告预测误差、并列集合式 top-$k$ 命中、regret、精确恢复合法率及同预算在线质量--成本，按留出 case 汇总配对差值与不确定性。61 个 recovery 并列不任意挑选赢家。四通道优于三通道及 degree priority 有效均为待检验假设，不填入模拟增益。

\subsection{支持或否定核心假设的结果}

核心假设是“布局状态能够预测不同算法的相对短期收益”。

\begin{itemize}
    \item 若 2D CNN 相比全局指标 Bandit 在相同预算下提高 reward、降低 regret，支持空间上下文有用；
    \item 若加入 wiremask 后选择准确率与最终指标均改善，支持派生 mask 作为门控/约束，而不是装饰性模块；
    \item 若 3D/ConvLSTM 只提高预测精度，却不提高在线决策的最终结果或成本效率，则不能声称状态流学习带来系统收益；
    \item 若 Bandit 与最强固定算法无显著差异，应将结论限定为“能识别某些状态下的算法偏好”，而不是宣称普遍优于所有单一算法。
\end{itemize}

\section{开源实现与复用边界}

下表中的地址来自项目本地归档的 README 与开源代码清单；正式实验必须在运行前记录仓库 commit/hash、许可证和依赖版本。下列“可复用”只表示适合作为实现参考或 baseline，不表示所有代码都可直接作为透明 worker。

\begin{longtable}{p{0.18\textwidth}p{0.27\textwidth}p{0.18\textwidth}p{0.27\textwidth}}
\caption{三类算法、物理评估器及相关代码地址}\label{tab:opensource}\\
\toprule
类别 & 项目/地址 & 在本方案中的用途 & 复用边界 \\
\midrule
\endfirsthead
\toprule
类别 & 项目/地址 & 在本方案中的用途 & 复用边界 \\
\midrule
\endhead
\bottomrule
\endfoot
模拟退火 & \urlbreak{https://github.com/bu-icsg/TAP-2.5D} & SA worker、给定布局评价、快速线长与热/路由调用参考 & 需要处理 HotSpot、CPLEX 和原始输入格式；路由线长和 HPWL 不混用 \\
解析法 & \urlbreak{https://github.com/Brilight/ATPlace_pub} & analytical baseline、Bookshelf/JSON wrapper、HPWL 与热感知参数参考 & PlaceFlow 核心/部分评估代码加密；透明实验应保留公开 wrapper，并单独实现可审计的 analytical worker \\
强化学习 & \urlbreak{https://github.com/weiweihook/RLPlanner} & RL worker、快速热模型接口和 PPO 配置参考 & 输入格式与 Bookshelf 不同，需要 canonical adapter；README 中的加速/误差为作者报告，需在本项目上重新校准 \\
热模型/布局掩码参考 & \urlbreak{https://github.com/xunk33333/MaskPlace2.5D} & 紧凑热模型、thermal data processing、mask 概念参考 & 核心 placement code 未公开，不能当作完整透明 placement worker \\
综合布局参考 & \urlbreak{https://github.com/MicroSystemsDesignLab/STAMP2.5D} & 2.5D 布局/热约束相关实现和对照参考 & 先核对与本项目 benchmark、许可证和依赖的兼容性 \\
热评估器 & TAP-2.5D 与 ATPlace\_pub 中的 HotSpot wrapper；HotSpot 官方源码需按所采用版本固定 & Fast/\allowbreak Medium/\allowbreak Exact thermal evaluator 的接口参考与最终验证 & 快速模型不能替代最终 HotSpot；记录版本、网格、边界条件和材料参数 \\
\end{longtable}

\subsection{状态表征与时序建模参考}

以下工作用于支持网络组件、状态流和注意力设计的动机，不直接证明其在本项目 2.5D 芯粒布局中有效。最终结论仍以本项目统一 evaluator、固定预算和跨实例消融为准。

\begin{itemize}
    \item CBAM：\urlbreak{https://openaccess.thecvf.com/content_ECCV_2018/html/Sanghyun_Woo_Convolutional_Block_Attention_ECCV_2018_paper.html}，通道/空间注意力的轻量设计参考；
    \item Swin Transformer：\urlbreak{https://arxiv.org/abs/2103.14030}，局部窗口和层次化视觉表示参考；
    \item PredRNN：\urlbreak{https://arxiv.org/abs/2103.09504}，空间--时间状态流建模参考；
    \item Self-Attention ConvLSTM：\urlbreak{https://ojs.aaai.org/index.php/AAAI/article/view/6819}，循环记忆与注意力结合的参考；
    \item TimeSformer：\urlbreak{https://arxiv.org/abs/2102.05095}，时空注意力分解参考；
    \item RL-DAS：\urlbreak{https://arxiv.org/abs/2403.02131}，布局设计空间中的状态、动作和长期优化背景参考。
\end{itemize}

\subsection{建议的代码目录}

为避免把第三方代码和新方法混成不可解释的整体，建议实验仓库采用以下分层：

\begin{verbatim}
experiment/
  canonical/          # Bookshelf/TAP/RLPlanner -> common IR
  evaluator/
    geometry.py       # deterministic legality
    wirelength.py     # HPWL / routed-WL adapter
    thermal.py        # compact / HotSpot / calibration
  workers/
    sa_tap.py
    analytical_atplace.py
    rl_rlplanner.py
  context/
    rasterizer.py     # X_t and wiremask
    cnn2d.py
    convlstm3d.py
  bandit/
    finite_arms.py
    linucb.py
    thompson.py
  logging/
    trajectory.py
    state_flow.py
    provenance.py
\end{verbatim}

第三方 worker 只负责产生候选布局；它们不允许自带另一套最终 reward。所有指标都经过 `evaluator/`，所有选择都经过 `bandit/`，这样才能证明提升来自状态条件化的算法组合，而不是来自某个基线内部的隐含规则。

\section{实施顺序与最小可行版本}

\begin{enumerate}
    \item \textbf{MVP-0：}只实现 canonical IR、GeometryEvaluator、HPWL、固定热模型和三种 worker 的离线运行接口。
    \item \textbf{MVP-1：}建立有限 arm 配置库，跑全臂采样，得到每个 case 的宏观状态转移数据。
    \item \textbf{MVP-2：}训练 2D CNN，先做收益预测，再接 LinUCB/Thompson Sampling；与固定算法、随机选择和全局标量 Bandit 对比。
    \item \textbf{MVP-3：}加入 wiremask 的派生门控与动作约束，验证它是否真正改变选择质量和局部修改效果。
    \item \textbf{MVP-4：}只有当静态上下文存在明确不足时，采集连续状态快照，先训练 CNN+GRU，报告历史状态带来的增益。
    \item \textbf{MVP-5：}若 GRU 证明空间转移仍有额外信息，再训练小型 ConvLSTM 或低分辨率局部注意力，并加入动作条件 flow critic。
    \item \textbf{MVP-6：}最后加入连续超参数优化或分层 Bandit；不得用它替代第一版有限配置库的归因实验。
\end{enumerate}

\section{预期论文叙事与风险边界}

论文应围绕一个可证伪的主张展开：

\begin{quote}
在统一的 2.5D 芯粒布局表示和物理评估预算下，空间化的布局状态与 wiremask 能够帮助 Contextual Bandit 选择更适合当前状态的布局优化器，从而在质量--成本权衡上优于固定算法和无状态组合策略。
\end{quote}

该主张对应四类证据：空间状态是否提高收益预测、Bandit 是否降低选择 regret、算法选择是否改善最终布局、改进是否在相同 evaluator 资源下成立。不要把“能够调用三个开源算法”写成核心创新，也不要把“3D CNN 能预测状态变化”直接写成“理解了物理规律”。

主要风险及处理方式如下：

\begin{itemize}
    \item \textbf{三种 worker 输入不一致：}先做 canonical adapter，统一坐标、单位和 pin offset。
    \item \textbf{快速热模型偏差：}用 HotSpot 做独立校准和最终验证，不让近似模型单独支撑最终结论。
    \item \textbf{解析法核心不透明：}将 ATPlace\_pub 作为外部 baseline，同时实现一个透明的 analytical worker 用于机制消融。
    \item \textbf{Bandit 学到 case ID 而非布局规律：}按实例划分数据，加入跨规模测试和 held-out topology。
    \item \textbf{3D 模型数据量不足：}先完成 2D 版本；3D 版本只使用固定长度的状态窗口，并与 2D 版本共享数据划分。
    \item \textbf{所谓多 agent 只是包装：}最终系统图必须标出唯一的决策变量 \(a_t\)、统一 evaluator、状态反馈和对照组，而不是只画多个模块名称。
\end{itemize}

\section{结论}

本方案建议将项目的核心从“多 agent 协同布局”收敛为“面向 2.5D 布局状态的状态化宏动作控制”。第一版先把算法选择和连续参数优化拆开，用有限配置库验证 CNN 上下文与 wiremask 是否有用；第二版先用 CNN+GRU 验证历史状态是否能带来额外决策收益，再按需使用 ConvLSTM 或低分辨率局部注意力建模空间状态流，最后以 flow critic 检验动作对整个优化流程的影响。这样既响应导师提出的“agent 自主选择算法和参数”，又保留 wiremask 对布局区域和动作约束的作用，并且每一项机制都能被明确的消融实验检验。

\section{三类 Agent 的最终定义与信息职责}

本项目的三类 Agent 定义为三种算法 Agent，而不是将 HPWL、thermal、order 三个指标分别包装成 Agent。\textbf{SA-Agent} 在 wiremask 和预算约束下执行局部交换、移动、旋转或扰动，并选择温度、降温率、邻域半径与接受准则；\textbf{Analytical-Agent} 执行解析式连续优化、合法化和局部修复，并选择步长、惩罚权重、平滑强度与迭代预算；\textbf{RL-Agent} 使用统一布局接口生成候选动作序列，并选择 rollout 长度、探索率、候选动作数和局部更新预算。

四张图是共享观察通道：priority、thermal proxy、HPWL 和 legality；PhysicsEvaluator 是唯一的物理裁判；Contextual Bandit/Coordinator 选择 $(agent, config, mask, macro\_budget)$。因此 Agent 的身份来自算法和可执行动作，图来自状态，评估器负责裁决，三者不会被混为一谈。

\section{四通道的形式化图表示（fourmap v2）}

\paragraph{连接优先级。} 对芯粒 $i$，$d_i$ 是关联 pin 的出现次数（incident-pin count），不使用 net weight、唯一邻居数、时序关键度或计划放置顺序。定义
\[
q_i=\frac{d_i}{\max(\max_j d_j,1)},\qquad
M_{\rm priority}(g)=
\begin{cases}
\max_{i:\,|g\cap R_i|>0}q_i,&\exists i:\ |g\cap R_i|>0,\\
-1,&\text{背景}.
\end{cases}
\]
同度保留并列，孤立芯粒为 0，与背景 $-1$ 区分。多个芯粒覆盖同一像素时取最大 $q_i$，与输入遍历顺序无关；不再使用最后写入覆盖。历史 order 名称仅用于读取旧图，不引入 precedence 或 rank 语义。

\paragraph{热代理与 HPWL。} 新投影复用历史投影的 thermal proxy 与 HPWL 两通道。thermal proxy 是当前功耗/快速热模型的空间代理，不是摄氏温度。HPWL 将网络包围盒的线长贡献投影到网格，概念式为
\[
M_{\rm hpwl}(g)=\sum_n\frac{w_n\ell_n}{|B_n|}\mathbf{1}[g\in B_n],
\]
其中 $\ell_n$ 为未加权 HPWL，$w_n$ 只计一次；离散包围盒、退化包围盒与归一化遵循所记录的 projector 版本。精确 HPWL 标量由 evaluator 独立计算，不把栅格求和当作精确证明。绘图色标不改变训练输入。

\paragraph{合法性诊断。} 设 $R_i$ 为芯粒矩形，$g$ 为板内像素单元，$N$ 为芯粒数。重叠项按精确两两交集面积累计：
\[
L_{\rm overlap}(g)=\frac{1}{\max(N-1,1)}
\sum_{i<j}\frac{|g\cap R_i\cap R_j|}{|g|}.
\]
这保留 pair multiplicity：三重重叠累加三个 pair，不用交集并集，也不把同像素中的互不相交芯粒误算为重叠。

板宽高为 $B_x,B_y$。令 $o_{i,s}$ 为芯粒越过边 $s$ 的距离除以该边法向轴板长，$c_{i,s}(g)$ 为芯粒沿边切向区间覆盖该像素切向区间的比例，$b_s(g)$ 表示该边最近的一像素边界带，则
\[
L_{\rm boundary}(g)=\frac1N\sum_i\sum_{s\in\{L,R,B,T\}}
 o_{i,s}\,c_{i,s}(g)\,b_s(g),\qquad
M_{\rm legality}=L_{\rm overlap}+L_{\rm boundary}.
\]
当前实现使用单位板坐标；切向区间完全落在板外时，将覆盖权重 1 投到切向中心裁剪后的最近边界像素，避免完全板外芯粒消失。各越界边分别累加，角点可接收多边贡献。两项均不截断，不保证落在 $[0,1]$。无冲突背景为 0。

\textbf{精确 GeometryEvaluator 是唯一合法性裁判。} 单独保存 exact valid、overlap 与 boundary 标量及其容差；栅格不能证明合法性。数值容差附近非零 map 不必意味着 evaluator 判为非法，不能用 map 是否为零覆盖 exact verdict。

\paragraph{当前证据与问题边界。} 当前已有 2322 条标签、129 个状态，其中 61 个 recovery 状态存在并列最优；报告集合式命中率与 regret，不将并列强制改为单标签。27 个 HotSpot 校准样本的留出 MAE 为 $58.91\,{}^\circ\mathrm{C}$，拟合为负斜率仿射关系；这不支持将当前 thermal proxy 当作摄氏温度或可靠温度校准。Case4 原始 37 板问题不可行，59 板替代版本是新问题，必须另列 case 身份，不能作为原问题已恢复合法的证据。以上属于既有阶段记录。当前四通道 API 与 129 个状态、150 个 checkpoint、190 幅图的同步更新属于表示和可视化实现；CNN 训练及选择收益尚无结果。

\subsection{四通道的具体计算、单位与图例}

图内显示名依次固定为 \code{Connectivity priority}、\code{Spatial thermal proxy}、\code{HPWL contribution (mm/pixel)}、\code{Overlap / boundary violation}，与下述计算语义一一对应。

四通道数组按 priority、thermal proxy、HPWL、legality 排列。快照从左到右、状态流从上到下均采用这个顺序。网格原点为左下角，横轴对应 $x$、纵轴对应 $y$；$0$--$256$ 是原始像素坐标。叠加的 $16\times16$ 宏网格只帮助读图，每个宏格仍含 $16\times16$ 原始像素。PNG 采用 nearest 显示和 300 DPI，另导出 PDF/SVG，不用插值制造空间细节。

\paragraph{Priority 的计算和表现。} $d_i$ 是 netlist 中芯粒 $i$ 的 pin 出现次数，同一芯粒出现在不同网络中会重复计数；当前不按 net weight 加权、不去重为邻居数。$q_i=d_i/\max(\max_j d_j,1)$，像素覆盖冲突取最大 $q_i$。色条固定在 $[0,1]$，16 级 viridis 由深紫、蓝绿到黄色表示低、中、高连接度；浅灰表示背景 $-1$，黑色矩形表示芯粒边界。断连芯粒为 $0$，与空白区分。同连接度保留同色，大片同色不能为了美观而人为打散。netlist 不变时 $q_i$ 不随迭代改变，只有颜色块的位置变化。网页提供逐芯粒的 $d_i,q_i$ 数值表。这是连接度启发式，不能解释为执行顺序、计划放置顺序或时序关键路径。

\paragraph{热代理的计算和表现。} 令 $P_i$ 为芯粒功耗（W），$A_i$ 为物理面积（mm$^2$），$c_i$ 为单位板坐标的中心，$s_{ia}$ 为沿轴 $a$ 的归一化尺寸；当前实现为
\begin{equation}
\begin{aligned}
\sigma_{ia}^2&=s_{ia}^2/12+0.08^2,\\
M_T(g)&=\sum_i\frac{P_i}{A_i}
\exp\!\left[-\frac12\sum_{a\in\{x,y\}}
\frac{(u_{ga}-c_{ia})^2}{\sigma_{ia}^2}\right].
\end{aligned}
\end{equation}
这是功率密度启发的高斯扩散场，缺少热阻、层叠材料与真实散热边界，不能换成摄氏温度。inferno 色图由黑紫低值到橙黄高值，使用 $\gamma=0.6$ 的 PowerNorm 增强低中值差异；色条保留原代理数值，标注 Spatial thermal proxy，不写 $^\circ$C/K。图在像素中心采样，标量峰值在芯粒中心采样，二者采样位置不同。正式 27 个 HotSpot block 样本的留出 MAE 为58.91$^\circ$C，当前校准不支持温降结论。

\paragraph{HPWL 的计算和表现。} 引脚坐标包含芯粒半尺寸与真实 percentage offset。每条网络的 $L_n$ 按 mm 计算：最大和最小引脚 $x$ 之差，加最大和最小 $y$ 之差；令离散包围盒像素数为 $K_n$，每个覆盖像素累加 $w_nL_n/K_n$。退化包围盒至少分配一个像素。当前真实 Case 适配器为 two-pin nets，隐含 $w_n=1$。色条为 mm/pixel，magma 亮处表示更高线长贡献，浅灰处没有可显示贡献。SymLogNorm 的线性区固定为同 Case 共同最大值的1\%，其他区域对数压缩；仅在显示时屏蔽机器精度残差。它是引脚包围盒贡献场，不是实际导线路径、布线拥塞或物理梯度。

\paragraph{Legality 的计算和表现。} 按前文公式先求精确矩形交集，再将交集面积按像素比例分配；三芯粒重叠保留三个 pair 的贡献。越界距离除以对应板长，乘切向覆盖比例后投影到一像素宽板边带，并除以 $N$。两项相加，非负但不截断到 $[0,1]$。快照用线性 Reds：$0$ 为白色，正值从浅红到深红；板内红块定位重叠，板边红带定位越界。色条为无量纲违反强度，不是合法概率，其总和不能当作精确重叠率。图题分别报告 overlap、boundary 和 valid；精确 evaluator 的几何与容差决定最终判决，map 为零不构成合法性证明。

\paragraph{统一色阶与差分。} 每个 Case 的全部算法与五帧共用对应通道的范围，避免逐帧归一化掩盖变化。四通道差分都定义为 $\Delta M=M_{after}-M_{before}$，用以零为中心的 coolwarm 色图：蓝色减少、红色增加、白色近似不变。priority 的差分主要表示位置迁移；单处红蓝变化不能替代全局目标与成本。快照第四通道的红色表示绝对违反强度，差分图的红色表示违反增加，二者图题和色条分别标识。CNN 使用原始 float32 张量及训练阶段固定的归一化，不使用网页颜色作为物理输入。

\section{一步、预算与算法影响的公平比较}

定义一个 step 为：一个 Agent 在固定宏动作预算 $macro\_budget$ 内优化并返回一个候选布局，记录前后状态、动作、资源和 reward。当前同状态实验固定每个动作65次目标评价（1次起始评价与64次候选评价），候选的接受或拒绝均占预算；共享 baseline、合法化及最终核验另计。SA 的一次邻域交换、解析法的一次梯度迭代和 RL 的一次 rollout 不直接当作等价宏 step。正式报告资源向量 $(N_{eval},T_{wall},M_{RAM},M_{VRAM})$，同时单列梯度与策略训练成本；当前证明的是等目标评价次数，不声称等时间或等总算力。历史不等预算的 smoke 记录不混入正式比较。

\section{可审计图表与同步规范}

历史三联图与 $[3,H,W]$、$[T,3,H,W]$ 数据保留，不能改标签后当作新状态。
当前独立的 \code{state\_projection.py} / \code{project\_state} 生成 $[4,H,W]$；新快照、状态流与效果差分分别按 $[4,H,W]$、$[T,4,H,W]$ 和 before/after 四通道记录。
状态版本为 \code{connectivity-priority-}\allowbreak\code{thermal-hpwl-legality-v2}（fourmap v2），固定通道顺序 priority、thermal\_proxy、hpwl、legality。
图表附带 case、step、agent、config、seed、evaluator、state version 和归一化方式。

实现、状态数据与图表使用同一 state version 关联；历史记录保留其原有版本，保证表示变更可追溯。

\section{实验方案：四通道更新后的短期冻结版}

\subsection{冻结目标与核心主线}

本节是正式实验前的短期冻结基线。除非出现实现 bug、统一评估协议无法运行、全臂采样证明定义不可识别，或消融结果提供明确反证，否则不修改 Agent 身份、共享 map 语义、macro-step 定义和公平预算。

本项目的核心不是把三个算法机械串成流水线，而是学习一个状态条件化的算法组合策略：
\[
\text{布局状态}\rightarrow\text{CNN 上下文}\rightarrow
\text{选择算法/配置/掩码/预算}\rightarrow
\text{统一 PhysicsEvaluator}\rightarrow\text{状态反馈}.
\]
CNN 或后续的 ConvLSTM/局部注意力负责状态表示、收益预测或状态转移预测；它们不替代物理评估器，也不直接近似物理梯度。

\subsection{三类 Agent 的固定身份}

三类 Agent 是三种算法 worker，而不是三个目标指标或三个 evaluator：

\begin{itemize}
    \item \textbf{SA-Agent：}在 wiremask 和预算约束下执行交换、移动、旋转和随机扰动，配置包括初温、降温率、邻域半径、接受准则和局部预算，重点处理局部残余冲突、非光滑约束和需要跳出局部结构的状态；
    \item \textbf{Analytical-Agent：}执行解析式连续优化、合法化和局部修复，配置包括步长、平滑强度、密度/热惩罚、迭代数和合法化预算，重点处理目标函数仍有明显下降空间的全局结构；
    \item \textbf{RL-Agent：}使用统一布局接口生成候选动作序列，配置包括 rollout 长度、候选动作数、探索率和局部更新预算，重点处理具有重复拓扑分布、需要长程动作序列或已有策略先验的状态。
\end{itemize}

上层 Contextual Bandit/Coordinator 选择 $(agent, config, mask, macro\_budget)$；PhysicsEvaluator 是唯一裁判；priority、thermal proxy、HPWL、legality 是四个共享观察通道。第一版先使用有限配置库，再研究连续超参数优化。

\subsection{阶段偏好只作为可检验假设}

已有研究支持不同算法的优势随条件变化，但没有给出一个适用于所有本项目 2.5D case 的固定阶段顺序。TILOS 的宏布局比较显示 RL、SA 和解析法在 proxy cost、最终线长、时序与稳定性上并不由同一方法同时占优；TAP-2.5D、RLPlanner 和 ATPlace2.5D 分别代表热感知 SA、RL 与解析式路线，但它们的输入、热模型和预算并不完全统一。因此本文只提出待验证假设：粗粒度高重叠阶段可能更适合解析法或 RL 初始化，局部残余冲突可能更适合 SA，中间阶段由状态和预算决定。最终学习并报告
\[
\operatorname{Preference}(agent\mid state\_cluster, phase),
\]
不预先宣称“前期 analytical、中期 RL、后期 SA”为事实。

\subsection{四张 map 的冻结语义}

基础输入固定为 $[M_{\rm priority},M_{\rm thermal\_proxy},M_{\rm hpwl},M_{\rm legality}]$，标量旁路为 $v_t$。priority 使用关联 pin 次数除以最大次数、并列保留、背景 $-1$、重叠像素取最大值；thermal proxy 不代表摄氏温度；HPWL 沿用历史投影；legality 使用前文的精确 pair-intersection coverage 加归一化边界投影，不截断。wiremask 为派生动作约束或可选门控，不是基础第五通道。历史三图保留原始语义与版本，不回写旧 evaluator 或 reward。

\subsection{一步、轨迹、状态流与物理评估器}

一个 macro-step 是：某 Agent 在固定宏动作预算内执行若干内部微步，产生候选布局；候选布局经过一次统一 PhysicsEvaluator；系统记录前后 map、动作、指标与资源。SA 的一次邻域交换、解析法的一次更新和 RL 的一次 environment action 都是 micro-step，不能直接视为等价的一步。

算法轨迹记录 worker 内部全部微步；Bandit 样本是 $(X_t,a_t,X_{t+1},\Delta metrics,cost)$；状态流是在固定 macro-step 时间尺度上串联的 $X_t\rightarrow X_{t+1}\rightarrow\cdots$。第一版保存宏观转移，后续才采集固定窗口状态流。

PhysicsEvaluator 由 GeometryEvaluator、WirelengthEvaluator、ThermalEvaluator 和 CostTracker 组成。所有 worker 只产生候选布局，不得自带另一套最终 reward。快速热模型用于在线筛选，HotSpot/等价求解器用于校准和最终验证；CNN 的预测不是物理真值。

\subsection{短期冻结的实验顺序}

\begin{enumerate}
    \item \textbf{E0：}建立 canonical layout IR，统一坐标、单位、pin offset、约束和 evaluator 输出；
    \item \textbf{E1：}建立有限 arm 库
    \[
    A=\{(k,c)\mid k\in\{\mathrm{SA},\mathrm{Analytical},\mathrm{RL}\},c\in C_k\},
    \]
    在相同初始布局、seed、fidelity 和预算下进行全臂采样；
    \item \textbf{E2：}训练轻量 2D CNN 预测各 arm 的短期收益，再接 LinUCB 或 Thompson Sampling；wiremask 可作为四通道派生门控和动作约束；
    \item \textbf{E3：}仅当 E2 显示静态状态不足时加入 CNN+GRU，验证历史是否提升在线选择；
    \item \textbf{E4：}仅当 E3 显示状态变化有额外信息时加入小型 ConvLSTM 或低分辨率局部注意力；
    \item \textbf{E5：}最后才加入连续超参数和动作条件 flow critic，并将搜索成本计入 reward。
\end{enumerate}

\subsection{公平预算和关键消融}

正式比较锁定 $(N_{eval},T_{wall},M_{RAM},M_{VRAM})$，并报告最终目标、HPWL、热代理与独立高保真峰值温度、合法率、evaluator 次数、墙钟时间、峰值内存、失败率、regret 和 arm 命中率。若算法需要预训练，分别报告包含训练成本和已训练 worker 部署成本的结果。

主对照包括固定 SA、固定 analytical、固定 RL、随机 arm、全局标量 Bandit、2D CNN、2D CNN+wiremask、phase-conditioned Bandit，以及后续的 CNN+GRU 和 ConvLSTM/局部注意力。必须验证空间 map 是否比标量更有预测力、wiremask 是否真正减少无效动作、状态流是否改善在线结果，以及连续调参收益是否抵消搜索成本。注意力权重只作可视化线索，不能直接作为物理因果解释。

\subsection{开源实现边界}

    SA 参考 \urlbreak{https://github.com/bu-icsg/TAP-2.5D}，解析式布局参考 \urlbreak{https://github.com/Brilight/ATPlace_pub}，RL/快速热模型参考 \urlbreak{https://github.com/weiweihook/RLPlanner}，掩码/热模型参考 \urlbreak{https://github.com/xunk33333/MaskPlace2.5D}，综合 2.5D 参考 \urlbreak{https://github.com/MicroSystemsDesignLab/STAMP2.5D}，热评估器参考 \urlbreak{https://github.com/uvahotspot/HotSpot}。运行前固定 commit/hash、许可证、依赖和输入适配边界，第三方 worker 的结果必须回到本项目统一 evaluator。

\subsection{冻结条件}

当前短期冻结以下内容：三类算法 Agent 的身份、fourmap v2 的四通道语义、macro-step 的决策粒度、有限 arm + 2D CNN + Contextual Bandit 的主路线，以及统一 evaluator 和资源预算。只有实现 bug、协议不可运行、定义不可识别或明确反证才能修改，并同步相应实现、状态版本和实验协议。

\end{document}
